\chapter{Auxiliary Verbs}

% sein (to be)
\begin{verbentry}{
  style = {nom},
  toc = {sein (to be)},
  name = {sein},
  english = {to be},
  type = {irregular, + Nominativ},
  auxiliary = {sein}
}
 
     \vspace{5pt}
    \hrule
    \vspace{5pt}
  \begin{tabular}{rl}
     \multicolumn{2}{l}{\textbf{PRÄSENS}}\\
    ich & \textbf{bin}\\
    du & \textbf{bist}\\
    er/sie/es & \textbf{ist}\\
    wir & \textbf{sind}\\
    ihr & \textbf{seid}\\
    sie/Sie & \textbf{sind}\\
  \end{tabular}
  \hfill
     \begin{tabular}{rl}
    \multicolumn{2}{l}{\textbf{PRÄTERITUM}}\\
    ich & \textbf{war}\\
    du & \textbf{warst}\\
    er/sie/es & \textbf{war}\\
    wir & \textbf{waren}\\
    ihr & \textbf{wart}\\
    sie/Sie & \textbf{waren}\\
  \end{tabular}
    \hfill
     \begin{tabular}{rl}
    \multicolumn{2}{l}{\textbf{KONJUNKTIV II}}\\
    ich & \textbf{wäre}\\
    du & \textbf{wärest}\\
    er/sie/es & \textbf{wäre}\\
    wir & \textbf{wären}\\
    ihr & \textbf{wäret}\\
    sie/Sie & \textbf{wären}\\
  \end{tabular}
  \hfill
  \begin{tabular}{rl}
     \multicolumn{2}{l}{\textbf{IMPERATIV}}\\
   & \textbf{sei (du)!}\\
   & \textbf{seien wir!}\\
   & \textbf{seid (ihr)!}\\
   & \textbf{seien Sie!}\\
   & \vspace{10pt}
  \end{tabular}

    \vspace{10pt}
 
    \begin{tabular}{ll}
     \multicolumn{2}{l}{\textbf{PERFEKT}}\\
   &  sein + \textbf{gewesen}\\
  \end{tabular}
\hfill
   \begin{tabular}{ll}
   
    \multicolumn{2}{l}{\textbf{FUTURE I}}\\
   & werden + \textbf{sein}\\
 
  \end{tabular}
\hfill
   \begin{tabular}{ll}
      \multicolumn{2}{l}{\textbf{PLUSQUAMPERFEKT}}\\
   & waren + \textbf{gewesen}\\
  \end{tabular}
  
  \vspace{10pt}

   \begin{tabular}{lll}
      \multicolumn{3}{l}{\textbf{MIT PRÄPOSITIONEN}}\\
& \textbf{sein in} + Dat & (to be located in)\\
& \textbf{sein auf} + Dat & (to be on / at)\\
&\textbf{sein bei} + Dat & (to be with /at)\\
&\textbf{sein von} + Dat & (to be from / made of)\\
&\textbf{sein aus} + Dat & (to be from / made of)\\
&\textbf{sein für} + Akk & (to be for / in favour of)\\
&\textbf{sein mit} + Dat & (to be with / equipped with)
  \end{tabular}
\hfill
\begin{tabular}{lll}
\multicolumn{3}{l}{\textbf{TRENNBARE VERBEN}}\\
& \textbf{---} & \\
\end{tabular}
  \hfill
  % \begin{tabular}{lll}
  %    \multicolumn{3}{l}{\textbf{TRENNBARE VERBEN}}\\
% \vspace{70pt}
 % \end{tabular}



  \vspace{-5pt}
\end{verbentry}

% haben (to have)
\begin{verbentry}{
  style = {acc},
  toc = {haben (to have)},
  name = {haben},
  english = {to have},
  type = {irregular, + Akkusativ},
  auxiliary = {haben}
}
 
     \vspace{5pt}
    \hrule
    \vspace{5pt}
  \begin{tabular}{rl}
     \multicolumn{2}{l}{\textbf{PRÄSENS} }\\
    ich & \textbf{habe}\\
    du & \textbf{hast}\\
    er/sie/es & \textbf{hat}\\
    wir & \textbf{haben}\\
    ihr & \textbf{habt}\\
    sie/Sie & \textbf{haben}\\
  \end{tabular}
  \hfill
     \begin{tabular}{rl}
    \multicolumn{2}{l}{\textbf{PRÄTERITUM} }\\
    ich & \textbf{hatte}\\
    du & \textbf{hattest}\\
    er/sie/es & \textbf{hatte}\\
    wir & \textbf{hatten}\\
    ihr & \textbf{hattet}\\
    sie/Sie & \textbf{hatten}\\
  \end{tabular}
  \hfill
       \begin{tabular}{rl}
    \multicolumn{2}{l}{\textbf{KONJUNKTIV II} }\\
    ich & \textbf{hätte}\\
    du & \textbf{hättest}\\
    er/sie/es & \textbf{hätte}\\
    wir & \textbf{hätten}\\
    ihr & \textbf{hättet}\\
    sie/Sie & \textbf{hätten}\\
  \end{tabular}
  \hfill
  \begin{tabular}{rl}
     \multicolumn{2}{l}{\textbf{IMPERATIV} }\\
   & \textbf{hab(e) (du)!}\\
   & \textbf{haben wir!}\\
   & \textbf{habt (ihr)!}\\
   & \textbf{haben Sie!}\\
   & \vspace{10pt}
  \end{tabular}

    \vspace{10pt}
 
    \begin{tabular}{rl}
     \multicolumn{2}{l}{\textbf{PERFEKT}}\\
    &haben + \textbf{gehabt}\\
  \end{tabular}
\hfill
   \begin{tabular}{rl}
   
    \multicolumn{2}{l}{\textbf{FUTURE I} }\\
    &werden + \textbf{haben}\\
 
  \end{tabular}
\hfill
   \begin{tabular}{rl}
      \multicolumn{2}{l}{\textbf{PLUSQUAMPERFEKT}}\\
    &hatten + \textbf{gehabt}\\
  \end{tabular}
  


  \vspace{10pt}

   \begin{tabular}{lll}
      \multicolumn{3}{l}{\textbf{MIT PRÄPOSITIONEN}}\\
& \textbf{haben an} + Dat & (to have a problem / interest /fault)\\
& \textbf{haben für} + Akk & (to have for / feel towards)\\
&\textbf{haben von} + Dat & (to have heard / get from)\\
&\textbf{haben auf} + Akk & (to be in the mood for)\\
&\textbf{haben vor} + Dat & (to plan / intend)
  \end{tabular}
\hfill
\begin{tabular}{lll}
\multicolumn{3}{l}{\textbf{TRENNBARE VERBEN}}\\
& \textbf{---} & \\
\end{tabular}
 % \hfill
  % \begin{tabular}{lll}
   %   \multicolumn{3}{l}{\textbf{TRENNBARE VERBEN}}\\
%\vspace{50pt}
 % \end{tabular}




  \vspace{-5pt}
\end{verbentry}

% werden (to become)
\begin{verbentry}{
  style = {acc},
  toc = {werden (to become)},
  name = {werden},
  english = {to become},
  type = {irregular},
  auxiliary = {sein}
}
 
     \vspace{5pt}
    \hrule
    \vspace{5pt}
  \begin{tabular}{rl}
     \multicolumn{2}{l}{\textbf{PRÄSENS} }\\
    ich & \textbf{werde}\\
    du & \textbf{wirst}\\
    er/sie/es & \textbf{wird}\\
    wir & \textbf{werden}\\
    ihr & \textbf{werdet}\\
    sie/Sie & \textbf{werden}\\
  \end{tabular}
  \hfill
     \begin{tabular}{rl}
    \multicolumn{2}{l}{\textbf{PRÄTERITUM} }\\
    ich & \textbf{wurde}\\
    du & \textbf{wurdest}\\
    er/sie/es & \textbf{wurde}\\
    wir & \textbf{wurden}\\
    ihr & \textbf{wurdet}\\
    sie/Sie & \textbf{wurden}\\
  \end{tabular}
    \hfill
     \begin{tabular}{rl}
    \multicolumn{2}{l}{\textbf{KONJUNKTIV II} }\\
    ich & \textbf{würde}\\
    du & \textbf{würdest}\\
    er/sie/es & \textbf{würde}\\
    wir & \textbf{würden}\\
    ihr & \textbf{würdet}\\
    sie/Sie & \textbf{würden}\\
  \end{tabular}
  \hfill
  \begin{tabular}{rl}
     \multicolumn{2}{l}{\textbf{IMPERATIV} }\\
   & \textbf{werde (du)!}\\
   & \textbf{werden wir!}\\
   & \textbf{werdet (ihr)!}\\
   & \textbf{werden Sie!}\\
   & \vspace{10pt}
  \end{tabular}

    \vspace{10pt}
 
    \begin{tabular}{rl}
     \multicolumn{2}{l}{\textbf{PERFEKT}}\\
   & sein + \textbf{geworden}\\
  \end{tabular}
\hfill
   \begin{tabular}{rl}
   
    \multicolumn{2}{l}{\textbf{FUTURE I} }\\
    &werden + \textbf{werden}\\
 
  \end{tabular}
\hfill
   \begin{tabular}{rl}
      \multicolumn{2}{l}{\textbf{PLUSQUAMPERFEKT}}\\
   & waren + \textbf{geworden}\\
  \end{tabular}
  
\vspace{10pt}

\begin{tabular}{lll}
\multicolumn{3}{l}{\textbf{MIT PRÄPOSITIONEN}}\\
& \textbf{---} & \\
\end{tabular}
\hfill
\begin{tabular}{lll}
\multicolumn{3}{l}{\textbf{TRENNBARE VERBEN}}\\
& \textbf{---} & \\
\end{tabular}

\vspace{-5pt}
\end{verbentry}

% lassen (to let)
\begin{verbentry}{
  style = {default},
  toc = {lassen (to let)},
  name = {lassen},
  english = {to let},
  type = {irregular},
  auxiliary = {haben}
}
 
     \vspace{5pt}
    \hrule
    \vspace{5pt}
  \begin{tabular}{rl}
     \multicolumn{2}{l}{\textbf{PRÄSENS} }\\
    ich & \textbf{lasse}\\
    du & \textbf{lässt}\\
    er/sie/es & \textbf{lässt}\\
    wir & \textbf{lassen}\\
    ihr & \textbf{lasst}\\
    sie/Sie & \textbf{lassen}\\
  \end{tabular}
  \hfill
     \begin{tabular}{rl}
    \multicolumn{2}{l}{\textbf{PRÄTERITUM} }\\
    ich & \textbf{ließ}\\
    du & \textbf{ließt}\\
    er/sie/es & \textbf{ließ}\\
    wir & \textbf{ließen}\\
    ihr & \textbf{ließt}\\
    sie/Sie & \textbf{ließen}\\
  \end{tabular}
     \hfill
     \begin{tabular}{rl}
    \multicolumn{2}{l}{\textbf{KONJUNKTIV II} }\\
    ich & \textbf{ließe}\\
    du & \textbf{ließest}\\
    er/sie/es & \textbf{ließe}\\
    wir & \textbf{ließen}\\
    ihr & \textbf{ließet}\\
    sie/Sie & \textbf{ließen}\\
  \end{tabular}
  \hfill
  \begin{tabular}{rl}
     \multicolumn{2}{l}{\textbf{IMPERATIV} }\\
   & \textbf{lass(e) (du)!}\\
   & \textbf{lassen wir!}\\
   & \textbf{lasst (ihr)!}\\
   & \textbf{lassen Sie!}\\
   & \vspace{10pt}
  \end{tabular}

    \vspace{10pt}
 
    \begin{tabular}{rl}
     \multicolumn{2}{l}{\textbf{PERFEKT}}\\
   & haben + \textbf{gelassen}\\
  \end{tabular}
\hfill
   \begin{tabular}{rl}
   
    \multicolumn{2}{l}{\textbf{FUTURE I} }\\
    &werden + \textbf{lassen}\\
 
  \end{tabular}
\hfill
   \begin{tabular}{rl}
      \multicolumn{2}{l}{\textbf{PLUSQUAMPERFEKT}}\\
   & hatten + \textbf{gelassen}\\
  \end{tabular}
  
\vspace{10pt}

\begin{tabular}{lll}
\multicolumn{3}{l}{\textbf{MIT PRÄPOSITIONEN}}\\
& \textbf{---} & \\
\end{tabular}
\hfill
\begin{tabular}{lll}
\multicolumn{3}{l}{\textbf{TRENNBARE VERBEN}}\\
& \textbf{---} & \\
\end{tabular}

\vspace{-5pt}
\end{verbentry}
